%%%%%%%%%%%%
%
% $Autor: Müller, Hinrichs $
% $Datum: 2026-06-20 10:37:32Z $
% $Pfad: UltraschallradarAutomatisierungstechnik/Manuals/Benutzerhandbuch/Chapters/de/Troubleshooting.tex $
% $Version: 1 $
% !TeX spellcheck = en_GB/de_DE
% !TeX encoding = utf8
% !TeX root = manual 
% !TeX TXS-program:bibliography = txs:///biber
%
%%%%%%%%%%%%

\chapter{Fehlersuche}

\begin{figure}[h]
	\centering
	\begin{tikzpicture}[
	node distance=1.5cm and 2cm,
	startstop/.style={rectangle, rounded corners, minimum width=3cm, minimum height=1cm, align=center, draw=black, fill=red!20, thick},
	process/.style={rectangle, minimum width=3cm, minimum height=1cm, align=center, draw=black, fill=blue!10, thick},
	decision/.style={diamond, aspect=2.5, minimum width=3cm, minimum height=1cm, align=center, draw=black, fill=yellow!20, thick},
	arrow/.style={thick, -Latex}
	]
	
	\node (start) [startstop] {\textbf{Störung / Fehlerbildschirm}};
	\node (dec1) [decision, below=of start] {Leuchtet die\\rote LED dauerhaft?};
	\node (proc1) [process, right=of dec1, xshift=1cm] {Verkabelung (Sensor)\\prüfen \& USB trennen};
	\node (proc2) [process, below=of dec1, yshift=-0.5cm] {Processing-Software\\neu starten};
	\node (proc3) [process, below=of proc1] {Hardware-Reset\\(Button am Arduino)};
	
	\draw [arrow] (start) -- (dec1);
	\draw [arrow] (dec1) -- node[anchor=south] {\textbf{Ja (Hardwarefehler)}} (proc1);
	\draw [arrow] (dec1) -- node[anchor=east] {\textbf{Nein (Softwarehänger)}} (proc2);
	\draw [arrow] (proc1) -- (proc3);
	
\end{tikzpicture}
	\caption{Flussdiagramm: Fehlermeldungen}
\end{figure}

\section{Software-Fehlermeldungen}

Sollte während des Selbsttests oder im laufenden Betrieb ein kritischer Hardwarefehler auftreten oder die Verbindung unterbrochen werden, erscheint eine Meldung. Alle motorischen Bewegungen des Servos werden sofort gestoppt.

\begin{itemize}
	\item \textbf{Visuelle Darstellung:} Die Benutzeroberfläche wechselt auf einen warnenden, roten Hintergrund mit Warnhinweisen. Das Radarbild wird komplett ausgeblendet.
	\item \textbf{Fehlermeldung:} \textit{,, SYSTEMFEHLER  - Sicherstellen das der Computer mit dem Radarsystem verbunden ist!  - Anlage überprüfen und neustarten.''}
	\item \textbf{Hardware-Indikator:} Am Gehäuse des Radarsystems leuchtet die rote LED dauerhaft auf.
\end{itemize}

\begin{figure}[htbp]
	\centering
	\includegraphics[width=0.8\textwidth]{General/BenutzeroberflaecheBilder/BenutzeroberflaecheSystemFehler}
	\caption{Benutzeroberfläche: Systemfehler}
	\label{fig: Benutzeroberfläche: System fährt hoch}
\end{figure}

\subsection{Handlungsanweisung bei Systemfehler}
	
\begin{enumerate}
	\item \textbf{Stromversorgung trennen:} Das Radarsystm stromlos machen (USB-Verbindung zum Computer trennen).
	%\item \textbf{Verkabelung prüfen:} Den korrekten Sitz der Kabel am HC-SR04 Ultraschallsensor (GND, VCC, Trigger, Echo) sowie den Spannungsteiler auf 3,3V überprüfen.
	\item \textbf{Sichtfeld kontrollieren:} Sicherstellen, dass das Sichtfeld des Ultraschallensors beim Starten nicht direkt auf 0\,cm blockiert ist.
	\item \textbf{Neustart:} Das USB-Kabel neu verbinden und die Benutzeroberfläche neu starten, um den Selbsttest erneut zu initiieren.
\end{enumerate}






